%% AI for All -- the full deck (one PDF).
%% Master file: preamble + title + framing, then \input the seven part
%% fragments, one per section of the source.
%% Build:  latexmk ai_for_all.tex   (lualatex, two passes; see .latexmkrc)
%%
%% SOURCE. The whole deck is Julia Schneider and Lena Kadriye Ziyal,
%% "We Need to Talk, AI -- A Comic Essay on Artificial Intelligence" (2019),
%% free read at weneedtotalk.ai under CC BY-NC-SA 4.0. The comic is 60 pages
%% and seven named sections, and the deck is seven parts, one per section:
%%
%%   Part 1  Intro            pages  6-10
%%   Part 2  Basics           pages 11-19
%%   Part 3  Examples         pages 20-27
%%   Part 4  Chances          pages 28-36
%%   Part 5  Risks            pages 37-43
%%   Part 6  Outlook          pages 44-49
%%   Part 7  Social Utopias   pages 50-53
%%
%% ONE PAGE, ONE FRAME, SIX CAPTIONS. The comic sets every topic page as six
%% panels in a 2x3 grid with one caption under each, and those six captions
%% are the page's whole argument. A comic page becomes a frame, the frame
%% title is that page's own heading, the body is all six captions in the
%% comic's own reading order, and the right third is the page's artwork.
%% Nothing is paraphrased: the essay argues in its own voice or not at all.
%%
%% Pages 6 and 51 to 53 are the exceptions and keep the single-passage
%% \comicquote treatment, because the comic does not set them as panels:
%% page 6 is one line across an empty spread, and Social Utopias is four
%% pages of running prose.
%%
%% The captions were pulled off the PDF by line height (the caption font sets
%% at ~13.6pt leading, which nothing else on a page does), then every page was
%% transcribed and adversarially re-checked against a render of that page,
%% because the filter also catches speech bubbles, chart axis labels and
%% diagram labels that happen to share the font.
%%
%% ARTWORK IS NOT DRAWN YET. Every frame carries the grey placeholder strip
%% instead. The \sidecaption on each is the art brief: which panel of which
%% page goes there. Swap placeholder_gray for the real asset and the caption
%% stays as it is. The comic is CC BY-NC-SA 4.0, so panels may be reproduced
%% in teaching material as long as the credit stays on the slide. The credit
%% is the grey "We Need to Talk, AI, page NN" line that \panels and
%% \comicquote print under every body of quoted text; it names the work and
%% the page, which is what the licence asks for.
%%
%% ATTRIBUTED ONCE PER SLIDE, NOT TWICE. Those frames deliberately carry NO
%% \slidesources marker. Every one of them is the comic, the attribution line
%% already says so, and a second rotated [1] up the page edge on all 43 of
%% them was the same claim twice on the same slide. \slidesources stays where
%% there is no attribution line to carry it: the agenda, the frame that
%% explains how to read the deck, and the book's own reading list.
\documentclass[aspectratio=169]{beamer}
\usetheme{tuftedark}
\usetikzlibrary{calc,fit,shapes.geometric,positioning,arrows.meta}

% Drop slide images in images/; include by bare name, e.g.
%   \sideimagecover[0.33]{p17_datacloud}
\graphicspath{{images/}}

\title[AI for All]{AI for All}
\subtitle{We Need to Talk, AI: a comic essay, carried into the lecture}
\author[Schutera]{Prof.\ Dr.-Ing.\ Mark Schutera}
\institute{\textcolor{DHBWred}{DHBW}\ Ravensburg}
\date{\today}
\setbookversion{\today}{7 parts}

% Right-edge vertical watermark: the sentence the whole essay turns on, and
% the instruction it leaves the reader with.
\renewcommand{\tdwatermark}{Let your voice be heard}

% ---------------------------------------------------------------- citations
% One global numbering for the deck, one shared References frame at the end.
% sorting=none keeps the numbers in citation order, so the sources on one slide
% compress to a tidy range, e.g. [1-3].
\usepackage[style=numeric-comp,sortcites,sorting=none,backend=biber]{biblatex}
\addbibresource{../artificialintelligencefoundations.bib}
\renewcommand*{\bibfont}{\footnotesize}
\setlength{\bibitemsep}{0.35em plus 0.15em}
% url and doi go; howpublished stays, which is where the Google course track
% and the comic's own free-read address live. Stripping those too would cost
% the deck the only machine-readable pointer back to its sources.
\AtEveryBibitem{\clearfield{url}\clearfield{doi}\clearfield{urldate}\clearfield{note}}
\usepackage{multicol}% two-column References frame
\usepackage{array}% >{\raggedright} column prefixes in the panel grid

\setbeamercolor{bibliography entry author}{fg=tddim}
\setbeamercolor{bibliography entry title}{fg=tdfg}
\setbeamercolor{bibliography entry location}{fg=tddim}
\setbeamercolor{bibliography entry note}{fg=tddim}
\setbeamercolor{bibliography item}{fg=tddim}
\DeclareFieldFormat{labelnumberwidth}{\textcolor{tddim}{#1}}

% The side-edge source marker takes cite keys and prints their numeric labels
% ([1-3]) in the same vertical spot, in place of hand-typed author-years.
% The marker sits 3mm off the right page edge, which the side-image strip
% covers, so grey stops reading there; on those frames it goes red.
\makeatletter
\renewcommand{\slidesources}[1]{%
  \iftd@sideimage\def\td@srccol{DHBWred}\else\def\td@srccol{tddim}\fi
  \begin{tikzpicture}[remember picture,overlay]
    \node[rotate=90,anchor=west,inner sep=0pt,text=\td@srccol,font=\scriptsize]
      at ([xshift=-3mm,yshift=13mm]current page.south east)
      {\cite{#1}};
  \end{tikzpicture}%
}

% \sidecaption shadowed for this deck, as in the Intelligence Engineering
% deck. The theme sets the caption horizontal, white and italic INSIDE the
% strip, which a comic panel swallows; here it reads up the strip's outer
% edge in typewriter grey, mirroring \slidesources on the opposite page edge.
% It carries the art brief while the strip is still a placeholder, and the
% panel reference once it is not. 13mm clears the footer band.
\renewcommand{\sidecaption}[1]{%
  \setlength{\td@stripw}{\td@sidefrac\paperwidth}%
  \def\td@westtest{west}%
  \begin{tikzpicture}[remember picture,overlay]
    \ifx\td@sideside\td@westtest
      \node[rotate=90,anchor=west,inner sep=0pt,text=tddim,
            font=\ttfamily\scriptsize]
        at ([xshift=\dimexpr\td@stripw+2mm\relax,yshift=13mm]current page.south west){#1};
    \else
      \node[rotate=90,anchor=west,inner sep=0pt,text=tddim,
            font=\ttfamily\scriptsize]
        at ([xshift=\dimexpr-\td@stripw-2mm\relax,yshift=13mm]current page.south east){#1};
    \fi
  \end{tikzpicture}%
}
\makeatother

% ------------------------------------------------------------- deck helpers

% \panel{art brief} -- the standard side strip for this deck: a 0.33 cover
% cropped image on the right edge, captioned up its outer edge. One call
% instead of two, so every frame reserves exactly the same third and a later
% sweep that swaps placeholders for panels touches one macro, not 45 frames.
% Takes the real asset as an optional first argument once it exists:
%   \panel{p17, the data word cloud}              placeholder
%   \panel[p17_datacloud]{p17, the data word cloud}   the panel itself
\newcommand{\panel}[2][placeholder_gray]{%
  \sideimagecover[0.33]{#1}%
  \sidecaption{#2}%
}

% \begin{comicquote}{17} ...passage... \end{comicquote}
% One passage lifted verbatim from the comic, and the whole body of the
% frame. A red opening quotation mark, the passage a step above body size in
% the page colour, and the page it came from in grey underneath. The
% quotation mark is oversized and hangs into the left margin: it is the only
% thing on the slide that says the words are not mine.
%
% It opens the \sidebody itself rather than asking each frame to wrap it. A
% quoted passage on this deck ALWAYS sits in the column beside the panel, so
% the two belong in one macro; leaving it to the frame means a forgotten
% wrapper sets the passage across the full page width and it runs silently
% under the image strip. The body is grabbed (+b) so the passage can be
% written as ordinary prose across several lines without \par bookkeeping.
%
% The lead-in space is what stands in for vertical centring: \sidebody is a
% top-aligned minipage of unset height, so a \vfill inside it has nothing to
% distribute and the passage would sit hard under the title.
\NewDocumentEnvironment{comicquote}{m +b}{%
  \begin{sidebody}%
    \vspace*{0.10\paperheight}%
    \noindent\hspace*{-0.14em}%
    {\color{DHBWred}\fontsize{34}{34}\selectfont\bfseries\textquotedblleft}%
    \par\vskip-0.45em
    {\large #2\par}%
    \vskip1.0em
    {\color{tddim}\footnotesize We Need to Talk, AI\enspace\textbar\enspace page~#1\par}%
  \end{sidebody}%
}{}

% \panels{13}{cap1}{cap2}{cap3}{cap4}{cap5}{cap6}
% A comic page's SIX panel captions, in the comic's own 2x3 reading order:
% row 1 left, row 1 right, row 2 left, row 2 right, row 3 left, row 3 right.
% The six captions are the page's whole argument, so the frame carries all
% six; the artwork on the strip carries the rest.
%
% Six mandatory arguments rather than a \panelcap list on purpose. A macro
% cannot hand an alignment tab to a tabular (TeX's alignment scanner wants the
% & token itself, not one that arrives by expansion), so a list form would
% have to drop the grid; and six fixed slots make a missing caption a compile
% error rather than a silently short page.
%
% No standalone quotation mark here, unlike \comicquote. Six numbered cells
% under a "We Need to Talk, AI, page NN" line already read as a reproduction
% of a page rather than as my own writing, and the oversized mark cost some
% 30pt of column height that the densest pages (42, 19, 45, 47) need: page 42
% carries 851 characters across its six panels and overflows the frame
% without them. \linespread is set for the same reason, and the columns are
% near-touching so a caption wraps as few times as possible.
\newcommand{\panels}[7]{%
  \begin{sidebody}%
    % Negative, on purpose. Beamer leaves some 32pt between the theme's title
    % rule and the first line of a \sidebody, which is comfortable for a
    % single quoted passage and is exactly what the dense pages cannot spare.
    % Pulling the grid up reclaims it without touching the title band.
    \vspace*{-0.055\paperheight}%
    \scriptsize\linespread{0.93}\selectfont
    \setlength{\tabcolsep}{0pt}%
    \renewcommand{\arraystretch}{1.0}%
    % Ragged right, not justified: at 45mm a justified column hyphenates
    % roughly every other line ("abil-ity", "counterfac-tuals"), which on a
    % projected slide reads as damage.
    \begin{tabular}{@{}>{\raggedright\arraybackslash}p{0.475\linewidth}@{\hspace{0.05\linewidth}}>{\raggedright\arraybackslash}p{0.475\linewidth}@{}}
      \tdpanel{1}{#2} & \tdpanel{2}{#3}\\[0.35em]
      \tdpanel{3}{#4} & \tdpanel{4}{#5}\\[0.35em]
      \tdpanel{5}{#6} & \tdpanel{6}{#7}
    \end{tabular}%
    \par\vskip0.45em
    {\color{tddim}\scriptsize We Need to Talk, AI\enspace\textbar\enspace page~#1\par}%
  \end{sidebody}%
}

% One cell of the grid: the panel's number in red, small, hanging at the start
% of the caption so the six read as a numbered sequence and a lecturer can
% point at "panel four" without counting.
\newcommand{\tdpanel}[2]{%
  {\color{DHBWred}\bfseries\tiny #1}\enspace{#2}%
}

% \gterm{...} -- one term inside a frame raised to the accent colour. Used on
% the framing slides only; a quoted passage is never recoloured, because a
% highlight inside a quotation is an edit.
\newcommand{\gterm}[1]{\textcolor{DHBWred}{#1}}

% Live QR code: \slideqr{caption}{url}, as in the Intelligence Engineering
% deck. The modules are white and sit straight on the dark page with no tile
% under them, which makes the page itself the quiet zone.
\usepackage{qrcode}
\newcommand{\slideqr}[2]{%
  \begin{tikzpicture}[remember picture,overlay]
    \node[anchor=north east, inner sep=0pt]
      (slideqr) at ([xshift=-9mm,yshift=-5mm]current page.north east)
      {\color{white}\qrcode[height=9mm,level=M]{#2}};
    \node[anchor=east, inner sep=0pt, text=tddim, align=right,
      font=\ttfamily\scriptsize]
      at ([xshift=-3mm]slideqr.west){#1};
  \end{tikzpicture}%
}

% ------------------------------------------------------------- footer nav
% The theme's footer is beamer's mini-frame navigation: one hollow dot per
% frame, grouped under its section name. It is sized for a deck of about eight
% sections; this deck runs seven, but at some fifty frames the dots still
% crowd the band, and the part you are in matters more here than the frame.
% So the nav is replaced, at deck level, with the same information compressed:
% the part you are in on the left, a progress track in the middle, the counter
% on the right. The theme is untouched and the other decks keep their dots.
% Drawn with \rule rather than tikz on purpose: a tikz coordinate resolves a
% bare length register through pgfmath, which hands back a unitless number that
% tikz then reads as centimetres, and the bar blows past the maximum dimension.
\newlength{\aiprogwd}
\newlength{\aiprogfill}
\setlength{\aiprogwd}{0.30\paperwidth}
% The fraction goes through pgfmath, which works in floating point. Doing it in
% TeX dimen arithmetic does not survive the first pass: before the aux file is
% written \inserttotalframenumber is stale, the \divide leaves \aiprogwd whole,
% and \aiprogwd times a frame number past about 120 exceeds the maximum TeX
% dimension. The clamp covers the same stale-aux case, where the current frame
% number can exceed the recorded total and the fill would run past its track.
\newcommand{\aiprogressbar}{%
  \ifnum\inserttotalframenumber>0\relax
    \ifnum\insertframenumber<\inserttotalframenumber\relax
      % The ratio is parenthesised so it is evaluated first. pgfmath keeps its
      % intermediates in a dimen register, so multiplying the track width by the
      % frame number before dividing overflows once past frame 120.
      \pgfmathsetlength{\aiprogfill}{\aiprogwd*(\insertframenumber/\inserttotalframenumber)}%
    \else
      \setlength{\aiprogfill}{\aiprogwd}%
    \fi
  \else
    \setlength{\aiprogfill}{0pt}%
  \fi
  \raisebox{0.6ex}{\makebox[\aiprogwd][l]{%
    \rlap{\textcolor{tdrule}{\rule{\aiprogwd}{1.1pt}}}%
    \textcolor{tdfg}{\rule{\aiprogfill}{1.1pt}}}}%
}
\setbeamertemplate{footline}{%
  \begin{beamercolorbox}[wd=\paperwidth,ht=6mm,dp=5mm,leftskip=6mm,rightskip=6mm]{section in head/foot}%
    {\usebeamerfont{section in head/foot}\color{tddim}\insertsectionhead}%
    \hfill
    \aiprogressbar
    \hfill
    {\usebeamerfont{section in head/foot}\color{tddim}%
      \insertframenumber\,/\,\inserttotalframenumber}%
  \end{beamercolorbox}%
}

% One divider per part, the theme's own section page. \sectionsubtitle is set
% in the part file just before its \section and carries the comic's own page
% range, so the divider says where in the book the part sits.
\AtBeginSection[]{\begin{frame}[plain]\sectionpage\end{frame}}

\begin{document}

\begin{frame}[plain]
  \titlepage
\end{frame}

% ---------------------------------------------------------------- agenda
% The lecture in one slide: the comic's own seven sections, in its own order,
% under its own page numbers, each labelled with the question it answers.
% Side strip like every other frame, so the deck's shape is set from frame two.
\begin{frame}{Seven parts, one comic}
  \panel{cover, the section spiral}
  \begin{sidebody}
    \footnotesize
    \renewcommand{\arraystretch}{1.25}%
    \begin{tabular}{@{}llp{0.46\linewidth}@{}}
      \toprule
      Part & \alert{Section} & \multicolumn{1}{@{}l@{}}{\alert{The question it answers}}\\
      \midrule
      \alertbf{I}   & Intro          & why talk about AI at all\\
      \alertbf{II}  & Basics         & what AI is, and what made it work\\
      \alertbf{III} & Examples       & where it already runs\\
      \alertbf{IV}  & Chances        & what it is good for\\
      \alertbf{V}   & Risks          & what it costs us\\
      \alertbf{VI}  & Outlook        & who gets to shape it\\
      \alertbf{VII} & Social Utopias & what we would want it for\\
      \bottomrule
    \end{tabular}
  \end{sidebody}
  \slidesources{Schneider2019}
\end{frame}

% The contract for the deck: every body slide is one quoted passage. Said once,
% here, so no later frame has to explain itself.
\begin{frame}{How to read this deck}
  \panel{page 5, the authors at the table}
  \begin{sidebody}
    \vspace*{0.05\paperheight}
    One comic page, one slide, \gterm{all six panel captions}. Nothing on
    the following slides is my wording. The frame title is the page's own
    heading, the six captions are lifted verbatim and numbered in the
    comic's own reading order, and the page number under them tells you
    where to read the rest.
    \vskip1.0em
    {\color{tddim}The comic is a free read at weneedtotalk.ai under
    CC BY-NC-SA 4.0. It is an essay, not a textbook: it argues. Argue back.}
  \end{sidebody}
  \slideqr{weneedtotalk.ai}{https://weneedtotalk.ai}
  \slidesources{Schneider2019}
\end{frame}

%% Part I -- Intro (comic pages 6 to 10).
%% The section that establishes why the book exists: the fear inherited from
%% fiction, the two authors and what each is after, and the turn from
%% replication to augmentation that the rest of the book runs on.
%%
%% ONE PANEL CAPTION, ONE SLIDE. The comic sets every topic page as six
%% panels in a 2x3 grid with one caption under each. Each caption takes a
%% slide of its own so the room reads one claim at a time; the six-dot row
%% under it marks which panel is on screen, and six consecutive slides
%% sharing a frame title are one comic page, read in the comic's own order.
%% Captions are verbatim, checked against a render of each page.
%% \setpagethesis opens each run with that page's one line, mine, not theirs.
%% Fragment only: no preamble, \input by ai_for_all.tex.

\sectionsubtitle{pages 6 to 10}
\section{Intro}

% Page 6 is the one topic page the comic does NOT set as six panels: a single
% line across an otherwise empty spread. So it keeps the single-passage
% \comicquote treatment while every page after it runs as six slides.
\begin{frame}{Knowing what future we would like to have}
  \panel{page 6, the opening spread}
  \begin{comicquote}{6}
    Knowing the future is impossible. But what we can do: Knowing what
    future we would like to have. And then work on it.
  \end{comicquote}
\end{frame}

% The page that sets the posture. Golem and Frankenstein establish that the
% fear is old, the rubber boots make it domestic, and Hawking turns it into a
% question with a deadline. Panels 5 and 6 are one Hawking quotation split
% across two panels.
\setpagethesis{The fear is older than the technology. Hawking only gave it a deadline.}
\begin{frame}{Our creatures have always turned against us}
  \panel{page 7, Golem, Frankenstein and the rubber boots}
  \panelslide{7}{1}{Our creatures have always turned against us. This applies to the Golem of Jewish mythology as well as to Frankenstein.}
\end{frame}
\begin{frame}{Our creatures have always turned against us}
  \panel{page 7, Golem, Frankenstein and the rubber boots}
  \panelslide{7}{2}{Neither magic nor natural sciences have ever helped to make artificial humans submissive and controllable.}
\end{frame}
\begin{frame}{Our creatures have always turned against us}
  \panel{page 7, Golem, Frankenstein and the rubber boots}
  \panelslide{7}{3}{Parents whose kids prefer to wear rubber boots instead of weather-proof sandals even when it's hot outside know that.}
\end{frame}
\begin{frame}{Our creatures have always turned against us}
  \panel{page 7, Golem, Frankenstein and the rubber boots}
  \panelslide{7}{4}{An AI Take-Over is a hypothetical scenario in which AI takes control of the planet -- taking it away from us.}
\end{frame}
\begin{frame}{Our creatures have always turned against us}
  \panel{page 7, Golem, Frankenstein and the rubber boots}
  \panelslide{7}{5}{In 2011, Stephen Hawking said that `Success in creating AI would be the biggest event in history.'}
\end{frame}
\begin{frame}{Our creatures have always turned against us}
  \panel{page 7, Golem, Frankenstein and the rubber boots}
  \panelslide{7}{6}{`Unfortunately, it might also be the last, unless we learn how to avoid the risks.' How can we make AI systems both safe and beneficial?}
\end{frame}

% Julia, the economist. Panel 3 is the bias claim the whole Risks section
% later rests on, and it is made as a compliment to data rather than an
% accusation.
\setpagethesis{Data will not lie to you strategically. It will still lie to you.}
\begin{frame}{Who is Julia}
  \panel{page 8, Julia, the digit wall and the 404}
  \panelslide{8}{1}{Good question. That's something I ask myself, now and then.}
\end{frame}
\begin{frame}{Who is Julia}
  \panel{page 8, Julia, the digit wall and the 404}
  \panelslide{8}{2}{I love making sense of our world. I love to discover patterns in data, with code.}
\end{frame}
\begin{frame}{Who is Julia}
  \panel{page 8, Julia, the digit wall and the 404}
  \panelslide{8}{3}{I love data's friendliness, their lack of strategic answering. If they are biased it's not their fault but ours.}
\end{frame}
\begin{frame}{Who is Julia}
  \panel{page 8, Julia, the digit wall and the 404}
  \panelslide{8}{4}{Our brains are prone to some errors. Overconfidence. And discrimination against people who seem to be different from us.}
\end{frame}
\begin{frame}{Who is Julia}
  \panel{page 8, Julia, the digit wall and the 404}
  \panelslide{8}{5}{Artificial Intelligence (AI) might have the potential to help us find answers to problems we are not able to solve.}
\end{frame}
\begin{frame}{Who is Julia}
  \panel{page 8, Julia, the digit wall and the 404}
  \panelslide{8}{6}{For that reason, we need to demystify AI. Therefore, I'm writing this comic. Let your voice be heard.}
\end{frame}

% Lena, the illustrator, and the answer to why this is a comic and not a
% paper: the subject is a mirror, so the medium is one too.
\setpagethesis{Ask what the way we draw AI says about what we expect from it.}
\begin{frame}{Who is Lena}
  \panel{page 9, Lena, the iceberg of visual codes}
  \panelslide{9}{1}{Not one thing in particular, I guess.}
\end{frame}
\begin{frame}{Who is Lena}
  \panel{page 9, Lena, the iceberg of visual codes}
  \panelslide{9}{2}{I love visual language. I love to encrypt thoughts into pictures by exploiting society's visual memory.}
\end{frame}
\begin{frame}{Who is Lena}
  \panel{page 9, Lena, the iceberg of visual codes}
  \panelslide{9}{3}{Visuality is one of the most vivid encryption systems I know. We are constantly increasing and renewing it collectively.}
\end{frame}
\begin{frame}{Who is Lena}
  \panel{page 9, Lena, the iceberg of visual codes}
  \panelslide{9}{4}{It's fun to mess around with visual codes by recomposing associations. Pictures start to shift meanings and interfere in society's processes.}
\end{frame}
\begin{frame}{Who is Lena}
  \panel{page 9, Lena, the iceberg of visual codes}
  \panelslide{9}{5}{AI is a field that gives us an intimate view on society's desires and fears.}
\end{frame}
\begin{frame}{Who is Lena}
  \panel{page 9, Lena, the iceberg of visual codes}
  \panelslide{9}{6}{Looking from different angles might help us to overcome commonly known perspectives. Let your view be seen!}
\end{frame}

% Maschinen-Maria 1927, R2D2 1977, T-800 1984, Samantha 2013, Tars 2014.
% Five panels of fiction and one panel that turns the parade into an argument.
\setpagethesis{Every AI in the parade is a character. Not one of them is a system.}
\begin{frame}{Lessons from AI in Fiction}
  \panel{page 10, the five fictional robots, 1927 to 2014}
  \panelslide{10}{1}{}
\end{frame}
\begin{frame}{Lessons from AI in Fiction}
  \panel{page 10, the five fictional robots, 1927 to 2014}
  \panelslide{10}{2}{}
\end{frame}
\begin{frame}{Lessons from AI in Fiction}
  \panel{page 10, the five fictional robots, 1927 to 2014}
  \panelslide{10}{3}{}
\end{frame}
\begin{frame}{Lessons from AI in Fiction}
  \panel{page 10, the five fictional robots, 1927 to 2014}
  \panelslide{10}{4}{}
\end{frame}
\begin{frame}{Lessons from AI in Fiction}
  \panel{page 10, the five fictional robots, 1927 to 2014}
  \panelslide{10}{5}{}
\end{frame}
\begin{frame}{Lessons from AI in Fiction}
  \panel{page 10, the five fictional robots, 1927 to 2014}
  \panelslide{10}{6}{Seems, like we should not seek to replicate but to augment ourselves.}
\end{frame}

%% Part II -- Basics (comic pages 11 to 19).
%% The three ingredients the book names for the current wave: computer power,
%% better algorithms, data. Then the reminder that what came out of them is
%% narrow, and what a second AI Winter would cost.
%% Fragment only: no preamble, \input by ai_for_all.tex.

\sectionsubtitle{pages 11 to 19}
\section{Basics}

% Page 13. The book refuses the word "intelligence" for what is shipping, and
% names what it would take to earn it. The footnote on the same page concedes
% it will keep saying AI anyway, which is worth reading aloud.
\begin{frame}{What is Artificial Intelligence?}
  \panel{page 13, the essence of a new situation}
  \begin{comicquote}{13}
    Up to now, it's more adequate to speak of `artificially augmented
    intelligence'.
  \end{comicquote}
  \slidesources{Schneider2019}
\end{frame}

% Page 14, first of the three key indicators. The number is the point: not
% Moore's 18 months, but 3,5, and the page prints both so the gap shows.
\begin{frame}{Computer Power}
  \panel{page 14, the cost and power curve}
  \begin{comicquote}{14}
    The amount of computer power used for training AI has gone up to 3,5
    month-doubling time.
  \end{comicquote}
  \slidesources{Schneider2019}
\end{frame}

% Page 15, second indicator. One sentence for the whole shift from rule
% writing to learning, and it is the sentence the room already half knows.
\begin{frame}{Algorithms 1}
  \panel{page 15, a cat? a cat? a dog..?}
  \begin{comicquote}{15}
    A second key indicator are better algorithms. Now we can teach machines
    instead of programming them.
  \end{comicquote}
  \slidesources{Schneider2019}
\end{frame}

% Page 16 is the buzzword page: ML, DL, NN, and the AI bookshelf that holds
% more than the three. The quote is the one that separates DL from ML, which
% is the distinction the room most often gets wrong.
\begin{frame}{Algorithms 2}
  \panel{page 16, the AI bookshelf}
  \begin{comicquote}{16}
    DL uses artificial Neural Networks (NN) that learn patterns directly
    from the data they are fed.
  \end{comicquote}
  \slidesources{Schneider2019}
\end{frame}

% Page 17, third indicator, and the widest definition of data in the book.
% The page itself is a word cloud of everything that counts as a feature.
\begin{frame}{Data 1}
  \panel{page 17, the data word cloud}
  \begin{comicquote}{17}
    Any form of raw fact or figure is data. Whether on paper or in
    electronic form.
  \end{comicquote}
  \slidesources{Schneider2019}
\end{frame}

% Page 18 adapts Maslow into a data pyramid: collect, store, transform,
% explore, AI. Four of the five steps are data work, which is the whole
% argument of the page and the reason it is quoted rather than summarised.
\begin{frame}{Data 2}
  \panel{page 18, the data pyramid}
  \begin{comicquote}{18}
    AI is a journey that begins with data: collecting, storing, transforming
    and exploring it.
  \end{comicquote}
  \slidesources{Schneider2019}
\end{frame}

% Page 19 runs the two AI Winters and lands on what today's systems actually
% are. This is the page to hold the room on: everything in Chances and Risks
% is about narrow AI, not the thing the news means.
\begin{frame}{General AI}
  \panel{page 19, the two AI Winters}
  \begin{comicquote}{19}
    Despite their impressive progress and success, today's AI is narrow. Its
    tasks are often classification and need a lot of data and a lot of
    energy.
  \end{comicquote}
  \slidesources{Schneider2019}
\end{frame}

%% Part III -- Examples (comic pages 20 to 27).
%% Seven systems already running, ordered from the one nobody notices (Maps)
%% to the one strapped to a body (exoskeletons). Deliberately mundane: the
%% argument is that the deployment already happened.
%%
%% ONE PANEL CAPTION, ONE SLIDE. The comic sets every topic page as six
%% panels in a 2x3 grid with one caption under each. Each caption takes a
%% slide of its own so the room reads one claim at a time; the six-dot row
%% under it marks which panel is on screen, and six consecutive slides
%% sharing a frame title are one comic page, read in the comic's own order.
%% Captions are verbatim, checked against a render of each page.
%% \setpagethesis opens each run with that page's one line, mine, not theirs.
%% Fragment only: no preamble, \input by ai_for_all.tex.

\sectionsubtitle{pages 20 to 27}
\section{Examples}

% The first example is the one the room used this morning. Panels 1, 2 and 3
% are one sentence split across three panels by the authors' own ellipses.
\setpagethesis{You supplied the training data this morning, on the way in.}
\begin{frame}{Maps}
  \panel{page 21, two phones, the delta and the route}
  \panelslide{21}{1}{Using location data from smartphones\ldots{}}
\end{frame}
\begin{frame}{Maps}
  \panel{page 21, two phones, the delta and the route}
  \panelslide{21}{2}{\ldots{} the AI-application Maps can analyze how fast traffic is moving\ldots{}}
\end{frame}
\begin{frame}{Maps}
  \panel{page 21, two phones, the delta and the route}
  \panelslide{21}{3}{\ldots{} in real-time.}
\end{frame}
\begin{frame}{Maps}
  \panel{page 21, two phones, the delta and the route}
  \panelslide{21}{4}{Additionally, maps incorporate user-reported traffic incidents.}
\end{frame}
\begin{frame}{Maps}
  \panel{page 21, two phones, the delta and the route}
  \panelslide{21}{5}{The proprietary algorithms of Maps have access to vast amounts of data.}
\end{frame}
\begin{frame}{Maps}
  \panel{page 21, two phones, the delta and the route}
  \panelslide{21}{6}{This means, Maps can suggest the best and fastest routes for us.}
\end{frame}

% From recommending a product to generating one. Panels 5 and 6 are one
% sentence, and the second half is the part worth arguing about.
\setpagethesis{When the platform generates the product too, what is left to choose?}
\begin{frame}{Recommendation}
  \panel{page 22, the component tag cloud}
  \panelslide{22}{1}{All major internet companies -- or `platforms' -- refine their products and processes using AI.}
\end{frame}
\begin{frame}{Recommendation}
  \panel{page 22, the component tag cloud}
  \panelslide{22}{2}{If we're looking for books or shoes, internet shops use AI so that our purchase decision gets easier.}
\end{frame}
\begin{frame}{Recommendation}
  \panel{page 22, the component tag cloud}
  \panelslide{22}{3}{Likewise, streaming providers use AI to analyze our consumer habits down to the last detail.}
\end{frame}
\begin{frame}{Recommendation}
  \panel{page 22, the component tag cloud}
  \panelslide{22}{4}{We benefit from improved suggestions and rediscoveries of old pearls which we have missed so far.}
\end{frame}
\begin{frame}{Recommendation}
  \panel{page 22, the component tag cloud}
  \panelslide{22}{5}{Also, the platforms use the knowledge of our habits to create and offer new products similar to previous successes\ldots{}}
\end{frame}
\begin{frame}{Recommendation}
  \panel{page 22, the component tag cloud}
  \panelslide{22}{6}{\ldots{}or to let AI autonomously stir those components together that have worked great elsewhere.}
\end{frame}

% Ends on WeChat and the supervision of groups of twelve or more. Panel 1
% is the precondition for everything the rest of the page describes.
\setpagethesis{Invisible by design, which is the precondition for everything else here.}
\begin{frame}{Social Networks}
  \panel{page 23, what a wonderful match}
  \panelslide{23}{1}{Those of us who use social networks use AI frequently, without necessarily knowing it.}
\end{frame}
\begin{frame}{Social Networks}
  \panel{page 23, what a wonderful match}
  \panelslide{23}{2}{Social networks use AI to recommend friends, news, photos, tags, advertisements.}
\end{frame}
\begin{frame}{Social Networks}
  \panel{page 23, what a wonderful match}
  \panelslide{23}{3}{They even `curate' the posts of your friends according to general or personal interests -- with the help of AI.}
\end{frame}
\begin{frame}{Social Networks}
  \panel{page 23, what a wonderful match}
  \panelslide{23}{4}{Now imagine a network that knows even more than our habits, relations, thoughts, interests, likes and dislikes, our hobbies and whereabouts.}
\end{frame}
\begin{frame}{Social Networks}
  \panel{page 23, what a wonderful match}
  \panelslide{23}{5}{Our state of health or bank account for example.}
\end{frame}
\begin{frame}{Social Networks}
  \panel{page 23, what a wonderful match}
  \panelslide{23}{6}{In China, this already exists in the form of `WeChat'. Just recently, the Chinese government officially announced to supervise each group with 12+ members.}
\end{frame}

% Refuses the word autonomous for what ships today, then names the real
% blocker in panel 6: not control, explanation.
\setpagethesis{The blocker is not control. It is explanation.}
\begin{frame}{Self-driving Cars}
  \panel{page 24, the unforeseen circumstance}
  \panelslide{24}{1}{Today, we can already activate the autopilot function which automatically and comfortably drives our car or flies our plane.}
\end{frame}
\begin{frame}{Self-driving Cars}
  \panel{page 24, the unforeseen circumstance}
  \panelslide{24}{2}{But ultimately, we are responsible. The systems are not autonomous -- or self-driving.}
\end{frame}
\begin{frame}{Self-driving Cars}
  \panel{page 24, the unforeseen circumstance}
  \panelslide{24}{3}{Within a few years, this might change.}
\end{frame}
\begin{frame}{Self-driving Cars}
  \panel{page 24, the unforeseen circumstance}
  \panelslide{24}{4}{We would have plenty of room to sleep, work, read or play, sober or drunk after a party.}
\end{frame}
\begin{frame}{Self-driving Cars}
  \panel{page 24, the unforeseen circumstance}
  \panelslide{24}{5}{Or maybe not. Nearly every car accident involves some sort of unforeseen circumstances, the AI must confront for the first time.}
\end{frame}
\begin{frame}{Self-driving Cars}
  \panel{page 24, the unforeseen circumstance}
  \panelslide{24}{6}{Moreover, user tests show that we need autonomous systems to explain their decision to feel safe. Not easy.}
\end{frame}

% From the 2015 crossover to cancer detection to the military, then to
% adversarial examples. The Picasso panel next to the rifle is the one to show.
\setpagethesis{If one system reads a tumour and a target, who aims it?}
\begin{frame}{Identifying Images}
  \panel{page 25, rifle 2017 next to Picasso 1910}
  \panelslide{25}{1}{A lot of AI work deals with identifying patterns. At an enormous scale.}
\end{frame}
\begin{frame}{Identifying Images}
  \panel{page 25, rifle 2017 next to Picasso 1910}
  \panelslide{25}{2}{In 2015, computers made fewer mistakes than we do sorting images into predefined categories.}
\end{frame}
\begin{frame}{Identifying Images}
  \panel{page 25, rifle 2017 next to Picasso 1910}
  \panelslide{25}{3}{Since 2018, everybody can use a free tool helping AI beginners detect images.}
\end{frame}
\begin{frame}{Identifying Images}
  \panel{page 25, rifle 2017 next to Picasso 1910}
  \panelslide{25}{4}{While early image identification could be lifesaving\ldots{}}
\end{frame}
\begin{frame}{Identifying Images}
  \panel{page 25, rifle 2017 next to Picasso 1910}
  \panelslide{25}{5}{\ldots{}sometimes it is life-threatening. The military is a huge fan of it, too.}
\end{frame}
\begin{frame}{Identifying Images}
  \panel{page 25, rifle 2017 next to Picasso 1910}
  \panelslide{25}{6}{Unfortunately, AI is susceptible to hallucinations: caused by biased data or by hackers who, for example, make AI mistake guns for apes.}
\end{frame}

% Panel 1 carries the number that makes the case: if words are a tenth of
% the message, a text-only interface reads a tenth of it.
\setpagethesis{Words are a tenth of the message, so a text interface reads a tenth of you.}
\begin{frame}{Emotional AI}
  \panel{page 26, ten per cent of the message}
  \panelslide{26}{1}{Only $\sim$ 10\% of the emotional meaning of a message is conveyed through words.}
\end{frame}
\begin{frame}{Emotional AI}
  \panel{page 26, ten per cent of the message}
  \panelslide{26}{2}{To really understand us, you need our facial expressions, our tone of voice, our gestures plus our words.}
\end{frame}
\begin{frame}{Emotional AI}
  \panel{page 26, ten per cent of the message}
  \panelslide{26}{3}{As our relationship with AI is becoming more and more intimate\ldots{}}
\end{frame}
\begin{frame}{Emotional AI}
  \panel{page 26, ten per cent of the message}
  \panelslide{26}{4}{\ldots{} we need our AI devices (like assistant systems, cars, or robots) to sense and adapt to what we mean and not only to our words.}
\end{frame}
\begin{frame}{Emotional AI}
  \panel{page 26, ten per cent of the message}
  \panelslide{26}{5}{Indeed, Emotional AI could help us with some highly valuable tasks. Help autistic children, for example, to learn emotions.}
\end{frame}
\begin{frame}{Emotional AI}
  \panel{page 26, ten per cent of the message}
  \panelslide{26}{6}{Nonetheless, our freedom of thought might be at stake in hostile scenarios.}
\end{frame}

% The one example in the section whose argument is fit rather than scale.
% A body is a moving target, so the controller has to learn.
\setpagethesis{A body is a moving target, so the controller has to keep learning.}
\begin{frame}{Exoskeletons}
  \panel{page 27, elderly man or firefighter}
  \panelslide{27}{1}{Exoskeletons are an external framework we can wear to augment our natural physical ability and reduce strain and weaknesses.}
\end{frame}
\begin{frame}{Exoskeletons}
  \panel{page 27, elderly man or firefighter}
  \panelslide{27}{2}{But every one of us is different.}
\end{frame}
\begin{frame}{Exoskeletons}
  \panel{page 27, elderly man or firefighter}
  \panelslide{27}{3}{We need assistance that is specifically fit for us. We need to tailor our devices at an individual level.}
\end{frame}
\begin{frame}{Exoskeletons}
  \panel{page 27, elderly man or firefighter}
  \panelslide{27}{4}{AI can use real-time measurements of our body signals, our breathing rate or hip extension for example.}
\end{frame}
\begin{frame}{Exoskeletons}
  \panel{page 27, elderly man or firefighter}
  \panelslide{27}{5}{This way, the AI can get our individual profile right. AI can make our exoskeletons fit our needs -- elderly man or firefighter.}
\end{frame}
\begin{frame}{Exoskeletons}
  \panel{page 27, elderly man or firefighter}
  \panelslide{27}{6}{Our bodies are different and constantly changing. Nowadays, the only feasible way of translating this to robotics is through the help of AI.}
\end{frame}

%% Part IV -- Chances (comic pages 28 to 36).
%% Eight arguments for the technology, made properly: the scale argument, the
%% accuracy argument, the bias argument turned around, and then the quiet one
%% on page 36 that a chance also costs something.
%%
%% ONE PANEL CAPTION, ONE SLIDE. The comic sets every topic page as six
%% panels in a 2x3 grid with one caption under each. Each caption takes a
%% slide of its own so the room reads one claim at a time; the six-dot row
%% under it marks which panel is on screen, and six consecutive slides
%% sharing a frame title are one comic page, read in the comic's own order.
%% Captions are verbatim, checked against a render of each page.
%% \setpagethesis opens each run with that page's one line, mine, not theirs.
%% Fragment only: no preamble, \input by ai_for_all.tex.

\sectionsubtitle{pages 28 to 36}
\section{Chances}

% The bottleneck is not storage, it is us. Everything after this page in
% the section follows from panel 3.
\setpagethesis{The bottleneck was never storage. It is us.}
\begin{frame}{Dealing with Big Data}
  \panel{hundreds of us against the volume}
  \panelslide{29}{1}{Our power to analyze the tons of data we are now able to store and to collect can be a bottleneck.}
\end{frame}
\begin{frame}{Dealing with Big Data}
  \panel{hundreds of us against the volume}
  \panelslide{29}{2}{AI helps us deal with large and complex datasets in ways we have never seen before.}
\end{frame}
\begin{frame}{Dealing with Big Data}
  \panel{hundreds of us against the volume}
  \panelslide{29}{3}{Even with hundreds of us analyzing patterns in big data, the sheer volume simply overwhelms our capacities.}
\end{frame}
\begin{frame}{Dealing with Big Data}
  \panel{hundreds of us against the volume}
  \panelslide{29}{4}{Even if we could analyze data, but it proves to be exhausting and tedious, we can use rule-based AI systems -- expert systems -- that do the job for us.}
\end{frame}
\begin{frame}{Dealing with Big Data}
  \panel{hundreds of us against the volume}
  \panelslide{29}{5}{Traditional algorithms are bound to follow the same logic over and over, without flexibility or learning new paths with new data.}
\end{frame}
\begin{frame}{Dealing with Big Data}
  \panel{hundreds of us against the volume}
  \panelslide{29}{6}{Now enter AI. Because AI systems get smarter as more data is given to them, they are well suited for patterns and anomalies over the time.}
\end{frame}

% The lawyer experiment: 20 lawyers, 5 contracts, 30 issues, 4 hours.
% Panel 5 holds all four numbers.
\setpagethesis{Ninety-five against eighty-five. What would you need to know to believe it?}
\begin{frame}{Efficiency}
  \panel{twenty lawyers against one AI}
  \panelslide{30}{1}{In many situations, AI is simply faster and more accurate than we are. Let's look at an example*.}
\end{frame}
\begin{frame}{Efficiency}
  \panel{twenty lawyers against one AI}
  \panelslide{30}{2}{A legal platform carried out a competition between 20 experienced lawyers against a trained AI.}
\end{frame}
\begin{frame}{Efficiency}
  \panel{twenty lawyers against one AI}
  \panelslide{30}{3}{Competitors had to review 5 legal contracts in 4 hours and identify 30 legal issues.}
\end{frame}
\begin{frame}{Efficiency}
  \panel{twenty lawyers against one AI}
  \panelslide{30}{4}{They were rated according to how accurately they identified each issue. Who won?}
\end{frame}
\begin{frame}{Efficiency}
  \panel{twenty lawyers against one AI}
  \panelslide{30}{5}{Humans achieved, on average, an 85\% accuracy rate, the AI 95\%. Humans took 92 minutes to complete the task, the AI\ldots{} 26 seconds.}
\end{frame}
\begin{frame}{Efficiency}
  \panel{twenty lawyers against one AI}
  \panelslide{30}{6}{The human lawyers who competed against the AI in the experiment said, the tasks were very similar to what lawyers do every day.}
\end{frame}

% The sharpest reasoning in the book. Panel 5 is the point: a shared
% distortion does not average out, it compounds.
\setpagethesis{A shared distortion does not average out. It compounds.}
\begin{frame}{Cognitive Biases}
  \panel{gender-bias guiding you through the world}
  \panelslide{31}{1}{Our brain, as a result of evolution, was not optimized for rational decision-making or perfect diagnoses.}
\end{frame}
\begin{frame}{Cognitive Biases}
  \panel{gender-bias guiding you through the world}
  \panelslide{31}{2}{Instead, we strive for a competitive degree of `fitness' in our specific environments; quite fast and often satisfying.}
\end{frame}
\begin{frame}{Cognitive Biases}
  \panel{gender-bias guiding you through the world}
  \panelslide{31}{3}{Our brain power is biologically limited; therefore, we work with heuristics, trust our gut feeling or store \& retrieve wrong evidence.}
\end{frame}
\begin{frame}{Cognitive Biases}
  \panel{gender-bias guiding you through the world}
  \panelslide{31}{4}{We call this cognitive bias. There are many of them and we share them to some degree.}
\end{frame}
\begin{frame}{Cognitive Biases}
  \panel{gender-bias guiding you through the world}
  \panelslide{31}{5}{And this is what makes our distortions have large-scale consequences. Were they random, they would cancel out.}
\end{frame}
\begin{frame}{Cognitive Biases}
  \panel{gender-bias guiding you through the world}
  \panelslide{31}{6}{Used wisely, AI could indeed help us improve our diagnoses and decisions.}
\end{frame}

% Gives the machine the generation and keeps the judgement. The last clause
% of panel 4 is the position, and it has not aged.
\setpagethesis{The machine generates. Judging stays with us, and judging is the part we enjoy.}
\begin{frame}{Creativity}
  \panel{just a generic style transfer}
  \panelslide{32}{1}{AI can save time otherwise wasted on stupid, tedious work for creativity. But it can do even more.}
\end{frame}
\begin{frame}{Creativity}
  \panel{just a generic style transfer}
  \panelslide{32}{2}{When an AI sorts our photos, it is easier to make a photo-calendar with our children (in 28 positions) for our parents as a present.}
\end{frame}
\begin{frame}{Creativity}
  \panel{just a generic style transfer}
  \panelslide{32}{3}{When an AI puts together a playlist of up-and-coming producers of house music, it may inspire our next DJ-set.}
\end{frame}
\begin{frame}{Creativity}
  \panel{just a generic style transfer}
  \panelslide{32}{4}{We can also get creative directly with the AI, compiling data and coding it, in a narrow area. But the AI has no judgement.}
\end{frame}
\begin{frame}{Creativity}
  \panel{just a generic style transfer}
  \panelslide{32}{5}{Judging is an ability we possess. And we love doing it. This is where AI together with us can achieve great things, with a human tweak.}
\end{frame}
\begin{frame}{Creativity}
  \panel{just a generic style transfer}
  \panelslide{32}{6}{The smaller the role of humans in the process of art-creation becomes, the more AI questions traditional concepts of art.}
\end{frame}

% The chance here is not the diagnosis, it is the time the diagnosis frees
% up. Panel 2 is unusually blunt about what that time currently goes on.
\setpagethesis{The chance is not the diagnosis. It is the time the diagnosis frees up.}
\begin{frame}{Connection}
  \panel{the doctor as overpaid bookkeeper}
  \panelslide{33}{1}{Most doctors became doctors because they wanted to connect with people in order to heal them.}
\end{frame}
\begin{frame}{Connection}
  \panel{the doctor as overpaid bookkeeper}
  \panelslide{33}{2}{Today, they often act like overpaid bookkeepers: taking in and spitting back data, prescribing drugs, adjusting doses, ordering tests.}
\end{frame}
\begin{frame}{Connection}
  \panel{the doctor as overpaid bookkeeper}
  \panelslide{33}{3}{And it's still only retrospective, error-prone knowledge.}
\end{frame}
\begin{frame}{Connection}
  \panel{the doctor as overpaid bookkeeper}
  \panelslide{33}{4}{But we cannot connect to doctors who struggle to keep up with messy data from disparate connections. And they cannot connect to us.}
\end{frame}
\begin{frame}{Connection}
  \panel{the doctor as overpaid bookkeeper}
  \panelslide{33}{5}{AI can support proactive, individualized strategies.}
\end{frame}
\begin{frame}{Connection}
  \panel{the doctor as overpaid bookkeeper}
  \panelslide{33}{6}{Our doctors could get to know us better, having the time to learn how a disease uniquely affects us and educate us in the right way.}
\end{frame}

% Robotics without reprogramming: demonstration replaces the rule set,
% which is page 15's move applied to a body.
\setpagethesis{Show the robot instead of programming it: page 15's move, applied to a body.}
\begin{frame}{New Opportunities}
  \panel{the sensor jacket and glove}
  \panelslide{34}{1}{AI allows for fascinating new opportunities. Teaching robots without a single line of coding, for example.}
\end{frame}
\begin{frame}{New Opportunities}
  \panel{the sensor jacket and glove}
  \panelslide{34}{2}{Robots solve complex tasks today -- but these tasks are deterministic. At any point of time, you know in advance where your robot needs to be.}
\end{frame}
\begin{frame}{New Opportunities}
  \panel{the sensor jacket and glove}
  \panelslide{34}{3}{Every time a task or location changes, the robot's software must be re-programmed. This is time-consuming, expensive and difficult.}
\end{frame}
\begin{frame}{New Opportunities}
  \panel{the sensor jacket and glove}
  \panelslide{34}{4}{Even for experienced programmers, it is hard to come up with a set of verbal rules of how to move when movements and settings are tricky or unpredictable.}
\end{frame}
\begin{frame}{New Opportunities}
  \panel{the sensor jacket and glove}
  \panelslide{34}{5}{AI now enables us to teach robots how to move, either together with the robot, or, with jackets and gloves containing sensors.}
\end{frame}
\begin{frame}{New Opportunities}
  \panel{the sensor jacket and glove}
  \panelslide{34}{6}{From the recorded example-movements, automation scripts can be created and optimized -- without us having to write complex code. Fantastic!}
\end{frame}

% The least contested case in the book. Reading, hearing, typing, the home,
% the classroom and the ward, one per panel.
\setpagethesis{Six barriers, one per panel, and not one of them is about convenience.}
\begin{frame}{Inclusion}
  \panel{a wall missing several bricks}
  \panelslide{35}{1}{AI is also about helping us, about improving inclusion and our well-being, by removing barriers.}
\end{frame}
\begin{frame}{Inclusion}
  \panel{a wall missing several bricks}
  \panelslide{35}{2}{AI reads to us, it narrates the world around us and help us `see'. Speech-to-text AI helps us `hear', too.}
\end{frame}
\begin{frame}{Inclusion}
  \panel{a wall missing several bricks}
  \panelslide{35}{3}{AI helps us type (with words) or talks to us. Users of all ages are adopting virtual assistants because it is so simple to learn how to handle them.}
\end{frame}
\begin{frame}{Inclusion}
  \panel{a wall missing several bricks}
  \panelslide{35}{4}{Smart home devices help us do things easily around the house and live independently.}
\end{frame}
\begin{frame}{Inclusion}
  \panel{a wall missing several bricks}
  \panelslide{35}{5}{In the classroom, AI can identify individual challenges and offer more personalized approaches.}
\end{frame}
\begin{frame}{Inclusion}
  \panel{a wall missing several bricks}
  \panelslide{35}{6}{AI virtual nurses or therapists offer 24-hour-support, or someone to talk to in case we are shy.}
\end{frame}

% Closes the section by undercutting it. The chance is convenience and the
% bill is the analogue institution. Filed under Chances, not Risks, on purpose.
\setpagethesis{What did you stop leaving the house for, and what closed because of it?}
\begin{frame}{Comfort}
  \panel{mall\_blueprint.pdf}
  \panelslide{36}{1}{AI satisfies and reinforces our desire for comfort.}
\end{frame}
\begin{frame}{Comfort}
  \panel{mall\_blueprint.pdf}
  \panelslide{36}{2}{Almost all answers to our consumer, communication or information needs are just a mouse click or an app away. And most of the time, AI is involved.}
\end{frame}
\begin{frame}{Comfort}
  \panel{mall\_blueprint.pdf}
  \panelslide{36}{3}{It is tempting not to have to leave our couch at all, no matter what we want to do: eat, work, chat, watch a movie, listen to music\ldots{}}
\end{frame}
\begin{frame}{Comfort}
  \panel{mall\_blueprint.pdf}
  \panelslide{36}{4}{'Pre-AI'-institutions such as shops, local transport, restaurants, libraries, or doctors must react to our new convenience.}
\end{frame}
\begin{frame}{Comfort}
  \panel{mall\_blueprint.pdf}
  \panelslide{36}{5}{All institutions with predominantly analogue offerings are undergoing a tremendous process of change. But it might be worth it.}
\end{frame}
\begin{frame}{Comfort}
  \panel{mall\_blueprint.pdf}
  \panelslide{36}{6}{It is precisely these institutions that are social places where we can come together and exchange ideas. Where we can decelerate. `Waste' time.}
\end{frame}

%% Part V -- Risks (comic pages 37 to 43).
%% Five risks, and none of them is the take-over the Intro opened with. They
%% are surveillance, a fragmented public sphere, the labour market, data
%% security, and the energy bill. Each is already running.
%% Fragment only: no preamble, \input by ai_for_all.tex.

\sectionsubtitle{pages 37 to 43}
\section{Risks}

% Page 39, the through-wall radio work. The page walks from words and faces
% to gait to radio reflections, and the step that matters is the last one:
% identity without line of sight.
\begin{frame}{Surveillance}
  \panel{page 39, the stick figures behind the wall}
  \begin{comicquote}{39}
    That means the AI does not only know where we are. The AI can also
    distinguish what we are doing, whether we are standing, or eating.
  \end{comicquote}
  \slidesources{Schneider2019}
\end{frame}

% Page 40. The risk named is structural, not a bad actor: the public sphere
% has broken into pieces, and the pieces are privately owned.
\begin{frame}{Opinion-Forming and Media}
  \panel{page 40, the fragmented public sphere}
  \begin{comicquote}{40}
    AI influences public opinion-forming and the media. Our public sphere
    has become fragmented.
  \end{comicquote}
  \slidesources{Schneider2019}
\end{frame}

% Page 41 is the labour page, and it refuses the usual blue-collar framing:
% lawyers, medics and brokers are on the list next to drivers. The quote
% taken is the other half, what stays with us and why.
\begin{frame}{Future Work}
  \panel{page 41, the job ladder}
  \begin{comicquote}{41}
    We are better in understanding non-trivial, new situations, building
    relationships, creating context and meaning. Custom-fit solutions.
  \end{comicquote}
  \slidesources{Schneider2019}
\end{frame}

% Page 42 is the CCC programme: the data letter, penalties for outflows,
% declared update frequency, permanent state investment in secure open
% source, and education. The opening sentence is the one that earns them.
\begin{frame}{Data Security and Safety}
  \panel{page 42, the data letter inbox}
  \begin{comicquote}{42}
    Our blind trust in AI providers is misplaced. We all know that the
    amount of data that we are implicitly or explicitly forced to leave
    behind is too great.
  \end{comicquote}
  \slidesources{Schneider2019}
\end{frame}

% Page 43. The comparison is with a brain, not with a previous data centre,
% and the numbers on the page (33 zettabyte in 2018, 175 by 2025, 8\% of
% global energy by 2030) are the ones to read out next to the quote.
\begin{frame}{Energy Use}
  \panel{page 43, save energy, save the planet}
  \begin{comicquote}{43}
    But conventional hardware was not designed for AI, hence it needs much
    more computing power, much more energy, than our brains do.
  \end{comicquote}
  \slidesources{Schneider2019}
\end{frame}

%% Part VI -- Outlook (comic pages 44 to 49).
%% The section that stops describing and starts demanding: a device that
%% augments rather than entertains, women in the room, AI read as an economic
%% form rather than a technology, the ART framework, and the literacy without
%% which none of the rest can be argued about.
%%
%% ONE PANEL CAPTION, ONE SLIDE. The comic sets every topic page as six
%% panels in a 2x3 grid with one caption under each. Each caption takes a
%% slide of its own so the room reads one claim at a time; the six-dot row
%% under it marks which panel is on screen, and six consecutive slides
%% sharing a frame title are one comic page, read in the comic's own order.
%% Captions are verbatim, checked against a render of each page.
%% \setpagethesis opens each run with that page's one line, mine, not theirs.
%% Fragment only: no preamble, \input by ai_for_all.tex.

\sectionsubtitle{pages 44 to 49}
\section{Outlook}

% The budget follows the purchase decision and the purchase decision follows
% entertainment, so the device we got is the device we bought. Panel 4 is a
% product-management claim.
\setpagethesis{Your phone analyses you all day. What has it ever helped you think?}
\begin{frame}{Overall Human-Machine Intelligence}
  \panel{page 45, the smartphone as marketplace}
  \panelslide{45}{1}{We mainly interact with AI via our smartphone. A proprietary marketplace for countless uncurated applications (=platform economy).}
\end{frame}
\begin{frame}{Overall Human-Machine Intelligence}
  \panel{page 45, the smartphone as marketplace}
  \panelslide{45}{2}{On the smartphone, we use AI basically for entertainment (image recognition), communication (language and translation) and information (e.g. google).}
\end{frame}
\begin{frame}{Overall Human-Machine Intelligence}
  \panel{page 45, the smartphone as marketplace}
  \panelslide{45}{3}{Unbeknownst to us, AI improves performance and battery life, too. And: It analyzes us, its users.}
\end{frame}
\begin{frame}{Overall Human-Machine Intelligence}
  \panel{page 45, the smartphone as marketplace}
  \panelslide{45}{4}{Up to date, there is no smartphone that is especially good at helping us get better ideas or interpret information. There is no interface for that either.}
\end{frame}
\begin{frame}{Overall Human-Machine Intelligence}
  \panel{page 45, the smartphone as marketplace}
  \panelslide{45}{5}{The development budgets are driven by numbers of purchase. In turn, our purchase decision is based on entertainment and communication values, and the design of the device.}
\end{frame}
\begin{frame}{Overall Human-Machine Intelligence}
  \panel{page 45, the smartphone as marketplace}
  \panelslide{45}{6}{But couldn't we also develop a smartphone that makes us more creative, more intelligent? That understands like us? Really augment us?}
\end{frame}

% Opens on the human computers who ran the first calculations and closes on
% what a homogeneous development team reproduces. The parenthesis in panel 6 is
% the authors' own.
\setpagethesis{The first computers were women. The fear is not a takeover but a homogeneous room.}
\begin{frame}{Feminism}
  \panel{page 46, Ada Lovelace and Grace Hopper}
  \panelslide{46}{1}{The contribution of women to AI is usually underestimated, especially concerning the democratic applicability of computing.}
\end{frame}
\begin{frame}{Feminism}
  \panel{page 46, Ada Lovelace and Grace Hopper}
  \panelslide{46}{2}{They were the first computers, actually: running calculations in the first information network of the world.}
\end{frame}
\begin{frame}{Feminism}
  \panel{page 46, Ada Lovelace and Grace Hopper}
  \panelslide{46}{3}{What we fear is not an AI-takeover. We fear scientists' homogeneity leading them to (finally!) `create' intelligent beings -- in an egocentric way.}
\end{frame}
\begin{frame}{Feminism}
  \panel{page 46, Ada Lovelace and Grace Hopper}
  \panelslide{46}{4}{In an AI world, where `female' service robots are `hot' and submissive, how can children learn to treat women with respect, with dignity?}
\end{frame}
\begin{frame}{Feminism}
  \panel{page 46, Ada Lovelace and Grace Hopper}
  \panelslide{46}{5}{And many of those whose service jobs will be replaced belong to an already vulnerable group: minority women.}
\end{frame}
\begin{frame}{Feminism}
  \panel{page 46, Ada Lovelace and Grace Hopper}
  \panelslide{46}{6}{We need (different!) women to participate in the development and use of AI. Otherwise we'll just reproduce patriarchal structures.}
\end{frame}

% Reads AI through Weber's judge as automaton and Marx's worker as appendage
% of the machine, then names the ask. Credited on the page to Timo Daum.
\setpagethesis{The paradigm is not production and sale. It is collection and use.}
\begin{frame}{Digital Capitalism}
  \panel{page 47, Weber's judge and Marx's worker}
  \panelslide{47}{1}{Algorithms have always been an integral part of our modern capitalist societies.}
\end{frame}
\begin{frame}{Digital Capitalism}
  \panel{page 47, Weber's judge and Marx's worker}
  \panelslide{47}{2}{The idea of objectiveness, general abstract rules, quantification, de-personalization and the concept of equality are characteristics of modernity.}
\end{frame}
\begin{frame}{Digital Capitalism}
  \panel{page 47, Weber's judge and Marx's worker}
  \panelslide{47}{3}{Max Weber, e.g., pictures the judge as an automation who processes concrete cases (data) based on an algorithm, namely the law.}
\end{frame}
\begin{frame}{Digital Capitalism}
  \panel{page 47, Weber's judge and Marx's worker}
  \panelslide{47}{4}{Marx' workers are an appendage of the machine, alienated from the spiritual potencies of the capitalist working process.}
\end{frame}
\begin{frame}{Digital Capitalism}
  \panel{page 47, Weber's judge and Marx's worker}
  \panelslide{47}{5}{Digital capitalism is the ground on which AI thrives. The central paradigm is -- instead of production and sale -- the collection and use of data.}
\end{frame}
\begin{frame}{Digital Capitalism}
  \panel{page 47, Weber's judge and Marx's worker}
  \panelslide{47}{6}{Since today's AI is largely incomprehensible to us, we need to regulate and shape AI as a public service and snatch the omnipotence over AI from the few dominant, digital companies.}
\end{frame}

% After Virginia Dignum. Three pillars, one per definition panel, and panel 6
% refuses to let the acronym stand in for the work.
\setpagethesis{The acronym is the easy part. The governance it asks for does not exist.}
\begin{frame}{ART: Accountability, Responsibility, Transparency}
  \panel{page 48, the three equal pillars}
  \panelslide{48}{1}{Throughout the world, we share a growing awareness: We need to make AI safe, beneficial and fair for us.}
\end{frame}
\begin{frame}{ART: Accountability, Responsibility, Transparency}
  \panel{page 48, the three equal pillars}
  \panelslide{48}{2}{This requires the participation and commitment from many of us. It means training, regulation and awareness in 3 equal pillars:}
\end{frame}
\begin{frame}{ART: Accountability, Responsibility, Transparency}
  \panel{page 48, the three equal pillars}
  \panelslide{48}{3}{Accountability is the concept that AI should be held responsible for the results of its algorithms.}
\end{frame}
\begin{frame}{ART: Accountability, Responsibility, Transparency}
  \panel{page 48, the three equal pillars}
  \panelslide{48}{4}{This includes our responsibility to develop a framework for AI that incorporates our values.}
\end{frame}
\begin{frame}{ART: Accountability, Responsibility, Transparency}
  \panel{page 48, the three equal pillars}
  \panelslide{48}{5}{Transparency refers to the need to describe, inspect and reproduce the AI algorithms and results, and to manage the data used, in a fair way.}
\end{frame}
\begin{frame}{ART: Accountability, Responsibility, Transparency}
  \panel{page 48, the three equal pillars}
  \panelslide{48}{6}{But let's not fool ourselves. For this ART of AI, we need a new and more ambitious form of governance.}
\end{frame}

% Closes the Outlook on numeracy: none of the demands above can be argued by
% people whose brains signal pain at a number.
\setpagethesis{Who in this room went quiet when the first equation appeared?}
\begin{frame}{No Fear of Numbers}
  \panel{page 49, 1908 against MCMVIII}
  \panelslide{49}{1}{In order to understand how AI could improve our life, we propose a real revolution:}
\end{frame}
\begin{frame}{No Fear of Numbers}
  \panel{page 49, 1908 against MCMVIII}
  \panelslide{49}{2}{Healing the public shaming in our individual math-biographies that caused fear.}
\end{frame}
\begin{frame}{No Fear of Numbers}
  \panel{page 49, 1908 against MCMVIII}
  \panelslide{49}{3}{Many of our brains signal `pain' when we read numbers, think or talk about them.}
\end{frame}
\begin{frame}{No Fear of Numbers}
  \panel{page 49, 1908 against MCMVIII}
  \panelslide{49}{4}{Not only does this deprive life. Zero, e.g., is a fascinating tool allowing us to write `1908' instead of MCMVIII: 1,000+(1,000-100)+5+1+1+1.}
\end{frame}
\begin{frame}{No Fear of Numbers}
  \panel{page 49, 1908 against MCMVIII}
  \panelslide{49}{5}{Zero is also part of the binary code, fundamental to todays' computers and their language. Fundamental to AI, therefore.}
\end{frame}
\begin{frame}{No Fear of Numbers}
  \panel{page 49, 1908 against MCMVIII}
  \panelslide{49}{6}{With improved `math-esteem', more of us could and should discuss chances and risks of AI in our lives.}
\end{frame}

%% Part VII -- Social Utopias (comic pages 50 to 53).
%% The only section of the book set as running prose rather than panels: four
%% pages of essay that ask what we would point the technology at if we got to
%% choose. The frames follow the essay's own order, and the last one is the
%% question the book ends on.
%% Fragment only: no preamble, \input by ai_for_all.tex.

\sectionsubtitle{pages 50 to 53}
\section{Social Utopias}

% Page 51. The time comparison is the section's opening move and the best
% single antidote to the word revolution: spears for 300,000 years, song
% recommendations for under ten.
\begin{frame}{AI today is not Skynet}
  \panel{page 50, the section opener}
  \begin{comicquote}{51}
    For over 300,000 years we've been hunting with spears; for less than ten
    years we've been looking for our next song with the help of AI.
  \end{comicquote}
  \slidesources{Schneider2019}
\end{frame}

% Page 52, and the conditional the whole utopia hangs on. Quoted in full
% because the second half, reduce not increase, is the part that gets
% dropped when this argument is repeated from memory.
\begin{frame}{It depends on the training data}
  \panel{page 52, the application procedure}
  \begin{comicquote}{52}
    If we train AIs on discriminatory data, they will make discriminatory
    decisions and reduce, not increase, the opportunities, and resources of
    the underprivileged.
  \end{comicquote}
  \slidesources{Schneider2019}
\end{frame}

% Page 52 again, the motorway example: an AI that documents every
% inhabitant's involvement transparently and proposes solutions. The claim
% quoted is the modest one the authors actually make for it.
\begin{frame}{Raising the question, not settling it}
  \panel{page 52, the motorway through the city}
  \begin{comicquote}{52}
    An AI may not be able to resolve all injustices, but it can raise
    questions about how we want to live.
  \end{comicquote}
  \slidesources{Schneider2019}
\end{frame}

% Page 53. The authors turn on their own utopia in one paragraph, which is
% the move that makes the section worth teaching. Ask the room to answer the
% third question before moving on.
\begin{frame}{Or a creepy surveillance state}
  \panel{page 53, the same platform, read twice}
  \begin{comicquote}{53}
    Or is this vision making us create a creepy surveillance state? Why
    should the authorities always be well-meaning? How could we make sure
    they are?
  \end{comicquote}
  \slidesources{Schneider2019}
\end{frame}

% The book's last words, and the deck's assignment. Tbc. is the authors'.
\begin{frame}{What do you think?}
  \panel{page 53, Tbc.}
  \begin{comicquote}{53}
    But the first step has been made: What can AI help us with? Where do we
    have to be careful? What do you think? Let your voice be heard and your
    point of view be seen.
  \end{comicquote}
  \slidesources{Schneider2019}
\end{frame}

% Page 55, the book's own reading list, carried over whole. It is the shelf
% the essay was written off, and it is the only frame in the part that is not
% a quotation, so it takes a tabular rather than a comicquote.
\begin{frame}{Further reading, the book's own}
  \panel{page 55, the reading list}
  \begin{sidebody}
    \scriptsize
    \renewcommand{\arraystretch}{1.15}%
    \begin{tabular}{@{}p{0.30\linewidth}p{0.62\linewidth}@{}}
      \toprule
      \multicolumn{1}{@{}l}{\alert{Authors}} & \multicolumn{1}{l@{}}{\alert{Title}}\\
      \midrule
      Agrawal, Gans, Goldfarb & Prediction Machines: The Simple Economics of Artificial Intelligence, 2018\\
      Evans                   & Broad Band: The Untold Story of the Women Who Made the Internet, 2018\\
      Chollet, Allaire        & Deep Learning with R, 2018\\
      Lee                     & AI Superpowers: China, Silicon Valley, and the New World Order, 2018\\
      Bostrom                 & Superintelligence: Paths, Dangers, Strategies, 2014\\
      Zuboff                  & The Age of Surveillance Capitalism, 2018\\
      Russell, Norvig         & Artificial Intelligence: A Modern Approach, 2010\\
      Daum                    & Das Kapital sind wir: Zur Kritik der digitalen \"Okonomie, 2017\\
      \bottomrule
    \end{tabular}
  \end{sidebody}
  \slidesources{Schneider2019}
\end{frame}


% -------------------------------------------------- the technical track
% The comic answers "what is AI for us"; it does not teach anyone to build
% one. Google's foundational-courses track is where that half lives, and it
% is what the previous version of this deck was carried over from. Keeping
% the four courses on a frame of their own, with the hub address spelled out
% and a QR code next to it, is what stops that link from going missing when a
% deck is rewritten. \nocite puts all five entries in the References frame, so
% the addresses survive in the bibliography too.
\nocite{GoogleFoundational2025,GoogleIntroML2025,GoogleMLCC2025,
        GoogleProblemFraming2025,GoogleManagingML2025}
% Starred, so it relabels the footer for the closing frames without numbering
% itself into the seven parts or firing \AtBeginSection and adding a divider.
\section*{Sources}
\begin{frame}{If you want to build one}
  \panel{a laptop on the same table}
  \begin{sidebody}
    \vspace*{0.03\paperheight}
    This deck is about what AI does to us. For how it is built, work
    Google's \gterm{foundational courses} track, four courses in order:
    \vskip0.7em
    \footnotesize
    \renewcommand{\arraystretch}{1.25}%
    \begin{tabular}{@{}llp{0.30\linewidth}@{}}
      \toprule
      \alert{Course} & \alert{Length} & \multicolumn{1}{@{}l@{}}{\alert{Answers}}\\
      \midrule
      Introduction to Machine Learning & 20 min   & what ML is\\
      Machine Learning Crash Course    & 12 modules & how to build one\\
      Problem Framing                  & 45 min   & whether to at all\\
      Managing ML Projects             & 90 min   & how to keep it alive\\
      \bottomrule
    \end{tabular}
    \vskip0.9em
    {\color{tddim}\scriptsize
     developers.google.com/machine-learning/foundational-courses}
  \end{sidebody}
  \slideqr{the track}{https://developers.google.com/machine-learning/foundational-courses}
  \slidesources{GoogleFoundational2025}
\end{frame}

% -------------------------------------------------------- shared references
\begin{frame}{References}
  \begin{multicols}{2}
    \printbibliography[heading=none]
  \end{multicols}
\end{frame}

\begin{frame}[plain]
  \vfill
  \centering
  {\LARGE\bfseries Let your voice be \alertbf{heard}}\par\vskip0.4em
  {\LARGE\bfseries and your point of view be \alertbf{seen}}\par\vskip1.2em
  {\small Prof.\ Dr.-Ing.\ Mark Schutera \;\textbar\; DHBW Ravensburg}
  \vfill
\end{frame}

\end{document}
